Fix \(\gamma=(\Gamma_0,\Gamma_1,\Gamma_2)\in\mathfrak J_{\Delta,p}^{\mathrm{WH}}\), and choose an arbitrary admissible representative \((\beta,A,\mu)\) with population mean \(m\).  By \(\text{def:stable-nonzero-restricted-strip-covariance-image}\), the representative is stable, satisfies \(A\neq0\) and the restricted principal-log strip, and its covariances equal the supplied triple.  Therefore
\[
 \eta=(m,\Gamma_1,\Gamma_2)\in\mathfrak I_{\Delta,p}^{\mathrm{WH}}
\tag{1}
\]
by \(\text{def:stable-nonzero-restricted-strip-moment-image}\).

Apply \(\text{thm:stable-nonzero-two-lag-wiener-hopf-recovery}\) to (1).  The part of its proof through reconstruction of \(C_{+}^{\mathrm{WH}}\) uses only \((\Gamma_1,\Gamma_2)\), and establishes
\[
 r=\operatorname{rank}(A),\quad S=\operatorname{im}(A),\quad
 F=[T|_S]_O,\quad L^{\mathrm{WH}}=[B|_S]_O,
\tag{2}
\]
\[
 H^{\mathrm{WH}}=[J|_S]_O^{-1},\quad K^{\mathrm{WH}}=K,\quad
 C_{+}^{\mathrm{WH}}(t)=C_{+}(t)\quad(t>0).
\tag{3}
\]
That theorem also proves basis invariance of all quantities in (2)--(3).

The curve in (3) extends continuously to zero, so the finite-interval integral in \(\text{def:wiener-hopf-covariance-only-inverse}\) exists.  Substituting (3) into \(\text{lem:zero-lag-bin-covariance}\) gives
\[
 \begin{aligned}
 D^{\mathrm{WH},0}
 &=\Delta^{-1}\!\left[\Gamma_0-
 \int_0^\Delta(\Delta-t)
 \{C_{+}^{\mathrm{WH}}(t)+C_{+}^{\mathrm{WH}}(t)^{\mathsf T}\}\,dt\right]\\
 &=\Delta^{-1}\!\left[\Gamma_0-
 \int_0^\Delta(\Delta-t)
 \{C_{+}(t)+C_{+}(t)^{\mathsf T}\}\,dt\right]=D\succ0.
 \end{aligned}
\tag{4}
\]
Hence \(\lambda^{\mathrm{WH},0}=\operatorname{diag}(D)=\lambda\), and the mean identity in the same lemma yields
\[
 m^{\mathrm{WH},0}=\Delta\lambda=m.
\tag{5}
\]
The argument passed to \(\Phi_{\Delta,p}^{\mathrm{WH}}\) by the covariance-only definition is therefore exactly the element (1).  The finite-dimensional transform/root theorem applies and gives
\[
 \Psi_{\Delta,p}^{\mathrm{WH},0}(\gamma)
 =\Phi_{\Delta,p}^{\mathrm{WH}}(m,\Gamma_1,\Gamma_2)
 =(\beta,A,\mu).
\tag{6}
\]
Equations (2)--(6) prove totality of every covariance-only operation.  Their left sides depend only on \(\gamma\), the known positive bin width \(\Delta\), and an inessential basis.  Since the admissible representative was arbitrary, (6) also forces every representative to equal the same recovered triple, proving representative independence and exact identification from the three covariance matrices alone.